\section{Discussion}
\label{sec:discussion_future_work}

\subsection{Primary comparison with the fixed-schedule baseline}

Across the 22 post-selection evaluation seeds, the complete PD-PPO training
configuration achieved lower mean forecast loss and lower macro-averaged
normalized forecast loss than the validation-selected fixed-schedule baseline in
every seed. The separate 24-seed descriptive
aggregate, which includes two pilot/model-selection seeds, retains the same
direction. With the ridge forecaster, the paired difference against the fixed
baseline selected for that forecaster is positive in 21 of the 22 post-selection
seeds and 23 of 24 seeds in the descriptive aggregate. Evaluation over all
scoreable test steps also
retains lower values of both co-primary endpoints in every seed, so the observed
difference is not confined to the event-balanced evaluation windows.
Executed-action analysis shows that the scheduler activates the corresponding
specialist channel for each of the particle, flux, and thermal event types while
retaining a broader mixture during calm periods. This pattern was observed for
the complete training configuration, which includes simulated warning inputs and
future-condition supervision; it cannot be attributed to the scalar objective
alone.

The reported system has a mandatory meteorological core channel and an optional-channel
capacity of one. The underlying design problem extends to settings in which only
$q<n_{\mathrm{opt}}$ optional
channels can operate. In such systems, a scheduler must decide whether to
reserve the available capacity for one generally useful measurement or
reallocate it as forecast requirements change. The present $q=1$ optional-channel capacity makes that trade-off directly observable and
prevents simultaneous low-cost optional channels. The meteorological core channel
costs 0.25 normalized units and each event-specific channel costs 0.49. Although
two event-specific channels would exceed the 0.75 budget, such subsets are absent
from the candidate set. Every listed candidate has steady-state demand at most
0.74 and startup-peak demand at most 0.92, below the corresponding 0.75 and 0.95
limits. Both instantaneous power limits are enforced, but the feasible-action
mask therefore excludes no candidate in this instance. The empirical result is
established under the optional-channel cardinality constraint and minimum
dwell-time constraint, not under observed binding of the two power limits
(Supplementary Table S14).

\subsection{Condition-dependent activation relative to the fixed schedule}

Calibration/validation selects a strong fixed optional channel. The FC4 snow
mass-flux channel is selected in more than half of the seeds, and the infrared
snow-surface-temperature channel in most of the remainder. Such a schedule
avoids switching and can perform well when the evaluation distribution is
dominated by one measurement requirement. Its weakness appears when particle,
flux, and thermal event types all contribute to the forecasting objective. The
optional channel preferred for one event type is nearly absent from another, so
no single subset can preserve the three event-specific benefits. This pattern
instantiates the proposition's motivation, but the empirical table does not
verify all of its policy-by-condition assumptions.

The PD-PPO scheduler obtains the observed difference with sparse executed-action
changes. The mean normalized Hamming distance per transition is 0.00369,
consistent with the minimum dwell-time constraint and below the corresponding
rate for the one-step look-ahead reference under its one-step rule. Executed
changes mainly coincide with changes in operating condition, and no optional
channel dominates the sequence across event types. Nonzero
condition--executed-action mutual information and the behavior-complexity gate
distinguish these executed sequences from fixed-like or simple-cycle-like
schedules under the frozen thresholds. These quantities are descriptive and
include boundaries between the eight separately reset event-balanced evaluation
windows.

\subsection{Scalar-objective and training-configuration comparisons}

The matched AoI-objective and uncertainty-objective PPO configurations replace
forecast loss with AoI or a bounded diagonal uncertainty proxy while retaining
the shared PPO algorithm. \Cref{prop:non_equivalence} establishes a mathematical
distinction among the theoretical forecast-loss, AoI, and covariance-trace
objectives; the empirical uncertainty configuration uses a bounded diagonal
proxy rather than the full covariance trace. All three PPO configurations
receive the same condition-specific training weights, training-only supervised
target actions, and auxiliary future-condition classification labels. Each
scalar objective also changes the value target, advantage, PPO surrogate, and
advantage weights used in the continuing behavior cloning term. The feasible
action set is small, and four optional-channel choices have nonzero activation
frequency under the main policy. The 95\% percentile bootstrap confidence
intervals for the mean paired differences include zero for both matched
scalar-objective PPO configurations. These comparisons concern matched PPO
training configurations with different scalar objectives; they are not
reward-only ablations.

The Double DQN training configuration uses the same scheduler observation,
reward based on forecast loss, action set, and constraints. It has higher
macro-averaged loss in all 22 post-selection seeds and higher
mean forecast loss in 21 of 22 seeds than PD-PPO. Under the channel-use
acceptance criterion, Double DQN meets the criterion in 5 of 22 seeds;
no channel is deactivated before warm-up completes, as structurally required by
$D_{\min}=6$. The PD-PPO configuration also
contains behavior cloning pretraining, continuing advantage-weighted behavior
cloning, and an auxiliary future-condition classifier. The comparison therefore
concerns complete training configurations and does not establish a general
advantage of policy-gradient learning over value-based learning.

The privileged one-step look-ahead reference using next-step test targets can
inspect next-step forecast loss on the test partition and therefore has more
information per decision than the deployed scheduler. It still has higher loss
in every seed and a mean normalized Hamming distance about 7.6 times as large.
The minimum dwell-time constraint, warm-up delay, and persistence of stale
observations couple consecutive actions, but the reference objective omits
their longer-term consequences. The result shows that next-step target access
under this one-step rule does not yield lower multi-step loss; it does not
isolate the value of sequential decision making or compare PD-PPO with a
constrained model-predictive controller.

\subsection{Effects of warning signals and true labels}

The warning-rule baseline and privileged true-label baseline test scheduler
behavior under different sources of condition-related information. The
simulated warning signals are leading scores generated from the simulated
operating-condition process. The warning-rule baseline applies the predeclared
threshold and a condition-to-action mapping selected on the
calibration/validation partition. The privileged true-label baseline removes
warning uncertainty by using the exact sequence $c_t,\ldots,c_{t+16}$. Neither
paired confidence interval excludes zero; the intervals establish neither a
population-level directional difference nor equivalence. The PD-PPO scheduler
instead combines simulated warning signals with partial observations, channel
sampling ages, channel operating states, and the previous executed channel
subset.

Both baselines depend on explicit condition-to-action mappings selected on the
calibration/validation partition. The true-label baseline uses condition labels
that are unavailable during deployment. The policy selects a proposed action
index from the feasible index set using current observations, channel operating states, and
recent-history features, and can retain information from observations between
warning transitions. Robustness to warning quality remains untested.
Over the continuous scoreable test partitions of the 22 post-selection seeds,
the thresholded four-class warning rule has 0.909 accuracy, a 0.916
macro-averaged F1 score, and 1.471 bits of mutual information with the current
condition label. Subtype-wise one-vs-rest ROC AUC is 0.997--0.999, and all 1,262
event onsets cross the 0.5 threshold within the preceding 16 steps, with median
lead time nine steps. These diagnostics establish that the simulator-generated
warning features are highly informative within the reported simulation; they do
not validate a field detector or test weak-warning conditions.
Supplementary Table S13 gives the component
metrics and analysis scope.

\subsection{Online computation and scaling}

The frozen primary forecaster is used for policy training and offline evaluation.
Online execution requires one policy forward pass
followed by a feasibility check and one check of the
minimum dwell-time constraint. It neither
simulates future observation sequences nor evaluates every candidate with a
forecaster. In the six-channel system, the categorical policy scores six
candidate indices; no runtime benchmark against instrument sampling and
estimation is reported.

With several simultaneously active optional channels, explicit enumeration can
grow combinatorially. A larger action space could retain the downstream
forecast-loss objective while requiring a different feasible-action
parameterization, such as sequential channel selection or direct feasible-subset
construction. The appropriate parameterization depends on the number of
channels and on whether several channel combinations carry distinct forecast
value.

\subsection{Limitations}

The study uses a controlled simulation for cold-region monitoring,
including simulated warning signals and six logical channels. The experiment was
designed with three event types, three corresponding specialist
channels, prescribed target weights, and leading warnings. It tests whether a
scheduler recovers a designed condition-dependent sensing allocation under the
specified information and execution structure. The 22 post-selection seeds and
two pilot seeds cover new event realizations from one generator, not different
sites, hardware assemblies, power systems, channel configurations, or warning
systems. Field use requires measured power and warm-up
profiles, warning signals derived from available instruments, and evaluation
under real missingness and faults. The privileged true-label baseline uses
evaluation-only information and is not part of the scheduler input.

The macro-averaged normalized forecast loss uses event-balanced evaluation
windows constructed from test truth using event coverage and transport intensity
so that all three event-specific measurement requirements receive explicit
support. Evaluation over all scoreable time steps in the test partition retains
the paired difference between the complete PD-PPO training configuration and the
validation-selected fixed-schedule baseline, but it does not replace evaluation
across different sites or event-generation processes outside the simulation
model. The reported result is established for mean forecast loss and
macro-averaged normalized forecast loss; invariance to alternative target or
event-type weightings is not claimed.

The primary comparison uses the frozen primary forecaster. The
alternative-forecaster analysis uses a separately fitted, frozen ridge
forecaster to evaluate existing observation sequences without policy retraining.
The scalar-objective comparisons also
retain behavior cloning pretraining, continuing advantage-weighted behavior
cloning, and auxiliary future-condition classification labels. The reported
implementation stores sampled proposed-action log probabilities even when the
previous action is executed, and the proposal-override rate was not logged; the
result therefore supports the complete implementation rather than an idealized
executed-action likelihood convention. Both power limits are non-excluding,
and the warm-up-interruption penalty is structurally zero in the reported
instance. Further work should test hardware, other channel sets, warning quality,
and joint scheduler--forecaster adaptation.
